\documentclass[aspectratio=169,11pt]{beamer}

% ------------------------------------------------
% PACKAGES - KEPT LIGHT FOR FAST COMPILATION
% ------------------------------------------------
\usepackage[T1]{fontenc}
\usepackage[utf8]{inputenc}
\usepackage{lmodern}
\usepackage{amsmath}
\usepackage{amssymb}
\usepackage{xcolor}

% ------------------------------------------------
% THEME
% ------------------------------------------------
\usetheme{Madrid}

% ------------------------------------------------
% COLORS
% ------------------------------------------------
\definecolor{Navy}{RGB}{25,45,75}
\definecolor{Teal}{RGB}{0,120,120}
\definecolor{Gold}{RGB}{205,145,35}
\definecolor{LightTeal}{RGB}{235,247,247}
\definecolor{LightGold}{RGB}{250,245,230}
\definecolor{DarkGray}{RGB}{55,55,55}

\setbeamercolor{structure}{fg=Teal}
\setbeamercolor{frametitle}{fg=white,bg=Navy}
\setbeamercolor{title}{fg=white,bg=Navy}
\setbeamercolor{author}{fg=white,bg=Navy}
\setbeamercolor{institute}{fg=white,bg=Navy}

\setbeamercolor{block title}{fg=white,bg=Teal}
\setbeamercolor{block body}{fg=black,bg=LightTeal}

\setbeamercolor{exampleblock title}{fg=white,bg=Gold}
\setbeamercolor{exampleblock body}{fg=black,bg=LightGold}

% ------------------------------------------------
% REMOVE NAVIGATION SYMBOLS
% ------------------------------------------------
\setbeamertemplate{navigation symbols}{}

% ------------------------------------------------
% FOOTER
% ------------------------------------------------
\setbeamertemplate{footline}{
  \leavevmode%
  \hbox{%
    \begin{beamercolorbox}[wd=.72\paperwidth,ht=2.5ex,
      dp=1ex,leftskip=1em]{author in head/foot}%
      Mathematics of Computing -- CUSAT M.Tech
    \end{beamercolorbox}%
    \begin{beamercolorbox}[wd=.28\paperwidth,ht=2.5ex,
      dp=1ex,rightskip=1em plus1fil]{date in head/foot}%
      \insertframenumber{} / \inserttotalframenumber
    \end{beamercolorbox}
  }%
}

% ------------------------------------------------
% TITLE INFORMATION
% ------------------------------------------------
\title[Mathematics of Computing]
{Mathematics of Computing}



\author{Sayona C Sajan}

\institute{
CUSAT -- M.Tech Artificial Intelligence \& Software Engineering
}

\date{2026}

% ------------------------------------------------
% DOCUMENT
% ------------------------------------------------
\begin{document}

% =================================================
% TITLE PAGE
% =================================================

{
\setbeamertemplate{footline}{}
\begin{frame}
    \titlepage
\end{frame}
}


% =================================================
% 1. PROOFS
% =================================================

\section{Proofs}

\begin{frame}{1. Proofs -- Questions}

\begin{enumerate}
\item \textbf{Computing:} A recursive program has
$T(n)=T(n-1)+1$, $T(1)=1$. Prove that $T(n)=n$.

\item Prove that the sum of two even integers is even.

\item Why can testing a finite number of cases never prove a universal mathematical statement?

\item What are the essential components of a mathematical proof?

\item Why is the argument ``all tested primes greater than 2 are odd, therefore every prime is odd'' invalid?
\end{enumerate}

\end{frame}

\begin{frame}{1. Proofs -- Answers}

\begin{enumerate}
\item \textbf{Induction:} Base case $T(1)=1$. Assume $T(k)=k$. Then
\[
T(k+1)=T(k)+1=k+1.
\]
Therefore $T(n)=n$.

\item Let $a=2m$ and $b=2n$. Then
\[
a+b=2(m+n),
\]
which is even.

\item Finite testing only establishes that the statement is true for the tested cases. A universal statement requires a proof for every allowed case.

\item Assumptions, definitions, logical reasoning, justified intermediate steps, and a conclusion.

\item The number $2$ is a counterexample. A universal statement fails if even one valid counterexample exists.
\end{enumerate}

\end{frame}

% =================================================
% 2. PROPOSITIONS
% =================================================

\section{Propositions}

\begin{frame}{2. Propositions -- Questions}

\begin{enumerate}
\item \textbf{Computing:} A system accepts a request only when the input is valid and the user is authenticated. Express this as a proposition.

\item What is a proposition?

\item When is an implication $p\to q$ false?

\item Express ``If the server is down, the application cannot connect'' symbolically.

\item Rewrite $p\to q$ using only $\neg$ and $\lor$.
\end{enumerate}

\end{frame}

\begin{frame}{2. Propositions -- Answers}

\begin{enumerate}
\item If $p$ = valid input and $q$ = authenticated user, then
\[
p\land q.
\]

\item A proposition is a declarative statement that is either true or false, but not both.

\item $p\to q$ is false only when
\[
p=T,\qquad q=F.
\]

\item If $p$ = server is down and $q$ = application can connect:
\[
p\to\neg q.
\]

\item
\[
p\to q\equiv\neg p\lor q.
\]
\end{enumerate}

\end{frame}

% =================================================
% 3. PREDICATES AND QUANTIFIERS
% =================================================

\section{Predicates and Quantifiers}

\begin{frame}{3. Predicates and Quantifiers -- Questions}

\begin{enumerate}
\item \textbf{Computing:} Express ``Every user has at least one valid login'' using quantifiers.

\item What is the negation of
\[
\forall x\,P(x)?
\]

\item Are $\forall x\exists y\,P(x,y)$ and
$\exists y\forall x\,P(x,y)$ equivalent?

\item Express ``There exists an integer whose square is 9.''

\item Negate
\[
\exists x\forall y\,P(x,y).
\]
\end{enumerate}

\end{frame}

\begin{frame}{3. Predicates and Quantifiers -- Answers}

\begin{enumerate}
\item
\[
\forall x\,
(User(x)\to\exists y\,ValidLogin(y,x)).
\]

\item
\[
\neg\forall xP(x)
\equiv
\exists x\neg P(x).
\]

\item No. The first allows a different $y$ for each $x$; the second requires one common $y$ for every $x$.

\item
\[
\exists x\in\mathbb Z\;(x^2=9).
\]

\item
\[
\forall x\exists y\,\neg P(x,y).
\]
\end{enumerate}

\end{frame}

% =================================================
% 4. TRUTH TABLES
% =================================================

\section{Truth Tables}

\begin{frame}{4. Truth Tables -- Questions}

\begin{enumerate}
\item \textbf{Computing:} A security rule is
\[
A=(p\land q)\lor\neg r.
\]
How many rows are required in its truth table?

\item When is an implication $p\to q$ true?

\item Prove using a truth table:
\[
\neg(p\land q)\equiv\neg p\lor\neg q.
\]

\item Define tautology, contradiction, and contingency.

\item Is $p\to q$ equivalent to $q\to p$?
\end{enumerate}

\end{frame}

\begin{frame}{4. Truth Tables -- Answers}

\begin{enumerate}
\item There are three variables. Therefore,
\[
2^3=8
\]
rows are required.

\item It is false only when $p$ is true and $q$ is false.

\item Evaluate both sides for all four combinations of $p,q$. Their truth values are identical, so they are logically equivalent.

\item A tautology is always true. A contradiction is always false. A contingency is sometimes true and sometimes false.

\item No. $p\to q$ and $q\to p$ are generally different propositions.
\end{enumerate}

\end{frame}

% =================================================
% 5. FOL SATISFIABILITY
% =================================================

\section{First-Order Logic Satisfiability}

\begin{frame}{5. First-Order Logic Satisfiability -- Questions}

\begin{enumerate}
\item \textbf{Computing:} Is
\[
\forall x(User(x)\to\exists y\,Login(y,x))
\]
satisfiable?

\item What does it mean for a formula to be satisfiable?

\item Is
\[
\forall xP(x)\land\exists x\neg P(x)
\]
satisfiable?

\item How can an existential formula be shown to be satisfiable?

\item Is every satisfiable formula valid?
\end{enumerate}

\end{frame}

\begin{frame}{5. First-Order Logic Satisfiability -- Answers}

\begin{enumerate}
\item Yes. Construct an interpretation where every user has at least one login.

\item A formula is satisfiable if there exists at least one interpretation in which it is true.

\item No. The first part says every object satisfies $P$, while the second requires one object that does not satisfy $P$.

\item Give a suitable domain and provide a witness satisfying the existential condition.

\item No. Satisfiable means true in at least one interpretation; valid means true in every interpretation.
\end{enumerate}

\end{frame}

% =================================================
% 6. PATTERN OF PROOF
% =================================================

\section{Pattern of Proof}

\begin{frame}{6. Pattern of Proof -- Questions}

\begin{enumerate}
\item \textbf{Computing:} A recursive algorithm changes $n$ to $n-1$. How can induction help prove its correctness?

\item How do you prove a universal statement?

\item How do you prove an existential statement?

\item How can a universal statement be disproved?

\item How does the form of a mathematical statement suggest the proof technique?
\end{enumerate}

\end{frame}

\begin{frame}{6. Pattern of Proof -- Answers}

\begin{enumerate}
\item Prove correctness for the base case, then assume correctness for $k$ and prove it for $k+1$.

\item Choose an arbitrary object and prove the property for it.

\item Provide a specific witness and prove that it satisfies the required property.

\item Give one valid counterexample.

\item Universal statements usually suggest direct proof or induction; existential statements suggest finding a witness; implications suggest proving the consequent from the antecedent.
\end{enumerate}

\end{frame}

% =================================================
% 7. PROOFS BY CASES
% =================================================

\section{Proofs by Cases}

\begin{frame}{7. Proofs by Cases -- Questions}

\begin{enumerate}
\item \textbf{Computing:} A program behaves differently for even and odd inputs. How can its correctness be proved?

\item Prove that $n(n+1)$ is always even.

\item Prove that $|x|\geq0$ using cases.

\item How could a proof involving an integer modulo 3 be divided into cases?

\item What is the most important requirement when using proof by cases?
\end{enumerate}

\end{frame}

\begin{frame}{7. Proofs by Cases -- Answers}

\begin{enumerate}
\item Consider the two exhaustive cases: $n$ even and $n$ odd. Prove correctness independently in each case.

\item If $n$ is even, the product is even. If $n$ is odd, $n+1$ is even. Hence the product is always even.

\item If $x\geq0$, then $|x|=x\geq0$. If $x<0$, then $|x|=-x>0$.

\item Consider:
\[
n\equiv0,1,2\pmod3.
\]

\item The cases must be exhaustive and should not leave out any possible situation.
\end{enumerate}

\end{frame}

% =================================================
% 8. IMPLICATION
% =================================================

\section{Proof of an Implication}

\begin{frame}{8. Proof of an Implication -- Questions}

\begin{enumerate}
\item \textbf{Computing:} A cache key being valid implies that the corresponding record exists. How can this implication be proved?

\item Prove: If an integer is divisible by 6, then it is divisible by 3.

\item What is the contrapositive of ``if $n^2$ is odd, then $n$ is odd''?

\item Is the converse of an implication automatically true?

\item What is vacuous truth?
\end{enumerate}

\end{frame}

\begin{frame}{8. Proof of an Implication -- Answers}

\begin{enumerate}
\item Assume the cache key is valid. Use the definition of validity to show that a corresponding stored record must exist.

\item Let $n=6k$. Then
\[
n=3(2k),
\]
so $n$ is divisible by 3.

\item The contrapositive is:
\[
n\text{ even}\to n^2\text{ even}.
\]

\item No. The converse requires a separate proof.

\item An implication is true whenever its antecedent is false, regardless of the consequent.
\end{enumerate}

\end{frame}

% =================================================
% 9. CONTRADICTION
% =================================================

\section{Proof by Contradiction}

\begin{frame}{9. Proof by Contradiction -- Questions}

\begin{enumerate}
\item \textbf{Computing:} Why can a graph path that visits more vertices than possible without repetition be used to derive a contradiction?

\item What is the basic method of proof by contradiction?

\item Give the main idea for proving that $\sqrt2$ is irrational.

\item Why does the statement ``there is no largest integer'' lead to contradiction if a largest integer is assumed?

\item When is contradiction especially useful?
\end{enumerate}

\end{frame}

\begin{frame}{9. Proof by Contradiction -- Answers}

\begin{enumerate}
\item A finite graph has only finitely many vertices. Assuming an excessively long simple path forces a repeated vertex, contradicting simplicity.

\item Assume the negation of the desired conclusion and derive an impossible statement.

\item Assume $\sqrt2=a/b$ in lowest terms. Squaring gives $a^2=2b^2$, forcing both $a$ and $b$ to be even, a contradiction.

\item If $n$ were largest, then $n+1$ would be a larger integer.

\item It is useful when directly proving the statement is difficult but its negation naturally creates an impossibility.
\end{enumerate}

\end{frame}

% =================================================
% 10. IFF
% =================================================

\section{Proving IFF}

\begin{frame}{10. Proving IFF -- Questions}

\begin{enumerate}
\item \textbf{Computing:} A binary tree is balanced iff the height difference of its subtrees is at most 1. How would you prove this?

\item What does $P\iff Q$ mean?

\item Prove:
\[
n\text{ is even}\iff n^2\text{ is even}.
\]

\item How can set equality be proved using iff?

\item Why is proving only one direction insufficient?
\end{enumerate}

\end{frame}

\begin{frame}{10. Proving IFF -- Answers}

\begin{enumerate}
\item Prove both directions: balanced $\to$ height difference at most 1, and height difference at most 1 $\to$ balanced.

\item It means both implications hold:
\[
P\to Q
\quad\text{and}\quad
Q\to P.
\]

\item If $n=2k$, then $n^2=4k^2$ is even. Conversely, if $n^2$ is even, $n$ must be even.

\item Show
\[
x\in A\iff x\in B
\]
for every element $x$.

\item One direction only establishes implication, not equivalence.
\end{enumerate}

\end{frame}

% =================================================
% 11. SETS
% =================================================

\section{Sets}

\begin{frame}{11. Sets -- Questions}

\begin{enumerate}
\item \textbf{Computing:} If $A$ is all users and $B$ is administrators, what does $A-B$ represent?

\item Define union and intersection.

\item What does $A\subseteq B$ mean?

\item How many subsets does a set containing $n$ elements have?

\item State De Morgan's laws for sets.
\end{enumerate}

\end{frame}

\begin{frame}{11. Sets -- Answers}

\begin{enumerate}
\item
\[
A-B
\]
represents users who are not administrators.

\item
\[
A\cup B
\]
contains elements in either set, while
\[
A\cap B
\]
contains elements common to both.

\item Every element of $A$ is also an element of $B$.

\item
\[
2^n.
\]

\item
\[
(A\cup B)^c=A^c\cap B^c
\]
and
\[
(A\cap B)^c=A^c\cup B^c.
\]
\end{enumerate}

\end{frame}

% =================================================
% 12. SET EQUATIONS
% =================================================

\section{Proving Set Equations}

\begin{frame}{12. Proving Set Equations -- Questions}

\begin{enumerate}
\item \textbf{Computing:} Show that the database query
``$(A$ OR $B)$ AND NOT $C$'' is equivalent to
\[
(A\land\neg C)\lor(B\land\neg C).
\]

\item Prove:
\[
A-B=A\cap B^c.
\]

\item Prove:
\[
A\cup(A\cap B)=A.
\]

\item What is the two-inclusion method?

\item Prove:
\[
A\cap(B\cup C)
=(A\cap B)\cup(A\cap C).
\]
\end{enumerate}

\end{frame}

\begin{frame}{12. Proving Set Equations -- Answers}

\begin{enumerate}
\item By distributivity:
\[
(A\lor B)\land\neg C
=
(A\land\neg C)\lor(B\land\neg C).
\]

\item
\[
x\in A-B
\iff x\in A\land x\notin B
\iff x\in A\cap B^c.
\]

\item By absorption:
\[
A\cup(A\cap B)=A.
\]

\item Prove both
\[
A\subseteq B
\]
and
\[
B\subseteq A.
\]
Then $A=B$.

\item Use distributivity:
\[
A\cap(B\cup C)
=(A\cap B)\cup(A\cap C).
\]
\end{enumerate}

\end{frame}

% =================================================
% 13. RUSSELL
% =================================================

\section{Russell's Paradox}

\begin{frame}{13. Russell's Paradox -- Questions}

\begin{enumerate}
\item \textbf{Computing:} Why does the collection of all data structures that do not contain themselves create a Russell-type paradox?

\item State Russell's paradox.

\item Why does unrestricted set formation cause problems?

\item How does axiomatic set theory avoid the paradox?

\item Why does
\[
R\in R\iff R\notin R
\]
produce a contradiction?
\end{enumerate}

\end{frame}

\begin{frame}{13. Russell's Paradox -- Answers}

\begin{enumerate}
\item Let $R$ contain all structures that do not contain themselves. Asking whether $R$ contains itself produces the contradiction.

\item Consider
\[
R=\{x:x\notin x\}.
\]
Then $R\in R$ iff $R\notin R$.

\item It allows definitions such as ``the set of all sets having property $P$'', which can be self-referential.

\item Modern axiomatic set theory restricts how sets can be constructed.

\item If $R\in R$, the definition says $R\notin R$. If $R\notin R$, the definition says $R\in R$. Both possibilities contradict themselves.
\end{enumerate}

\end{frame}

% =================================================
% 14. WELL ORDERING
% =================================================

\section{Well-Ordering Principle}

\begin{frame}{14. Well-Ordering Principle -- Questions}

\begin{enumerate}
\item \textbf{Computing:} Why does repeatedly selecting a smaller positive integer eventually terminate?

\item State the Well-Ordering Principle.

\item How is the minimal-counterexample method related to well-ordering?

\item Are all integers well-ordered under the usual $\leq$ relation?

\item What is the relationship between well-ordering and induction?
\end{enumerate}

\end{frame}

\begin{frame}{14. Well-Ordering Principle -- Answers}

\begin{enumerate}
\item The positive integers are well-ordered, so an infinite strictly decreasing sequence of positive integers cannot exist.

\item Every nonempty subset of the positive integers has a least element.

\item If counterexamples exist, choose the smallest one. Then derive a contradiction by constructing a smaller counterexample.

\item No. The set of all integers has no least element.

\item The Well-Ordering Principle and the Principle of Mathematical Induction are equivalent over the positive integers.
\end{enumerate}

\end{frame}

% =================================================
% 15. INDUCTION
% =================================================

\section{Mathematical Induction}

\begin{frame}{15. Induction -- Questions}

\begin{enumerate}
\item \textbf{Computing:} A loop processes one array element per iteration. How can induction prove that after $k$ iterations exactly $k$ elements have been processed?

\item State the Principle of Mathematical Induction.

\item Prove:
\[
1+2+\cdots+n=\frac{n(n+1)}2.
\]

\item What is the induction hypothesis?

\item What are common mistakes in induction proofs?
\end{enumerate}

\end{frame}

\begin{frame}{15. Induction -- Answers}

\begin{enumerate}
\item Prove the property for the first iteration, assume it after $k$ iterations, then show it after $k+1$ iterations.

\item Prove a base case and prove that
\[
P(k)\to P(k+1).
\]

\item Base:
\[
1=\frac{1(2)}2.
\]
Assume true for $k$. Then
\[
\frac{k(k+1)}2+(k+1)
=
\frac{(k+1)(k+2)}2.
\]

\item It is the assumption that $P(k)$ is true in the induction step.

\item Missing the base case, assuming the conclusion directly, or failing to prove the step from $k$ to $k+1$.
\end{enumerate}

\end{frame}

% =================================================
% 16. INVARIANTS
% =================================================

\section{Invariants}

\begin{frame}{16. Invariants -- Questions}

\begin{enumerate}
\item \textbf{Computing:} What loop invariant can be used for insertion sort?

\item What is a loop invariant?

\item What are the three major stages of using an invariant?

\item Why are invariants important in algorithm correctness?

\item Does proving an invariant alone prove that an algorithm terminates?
\end{enumerate}

\end{frame}

\begin{frame}{16. Invariants -- Answers}

\begin{enumerate}
\item Before each iteration, the processed portion of the array is sorted and contains exactly the same elements as the corresponding original portion.

\item A property that remains true before and after every loop iteration.

\item Initialization, maintenance, and termination.

\item They provide a structured way to prove that an algorithm maintains correctness throughout execution.

\item No. Termination must be established separately.
\end{enumerate}

\end{frame}

% =================================================
% 17. STRONG INDUCTION
% =================================================

\section{Strong Induction}

\begin{frame}{17. Strong Induction -- Questions}

\begin{enumerate}
\item \textbf{Computing:} For
\[
T(n)=T(n-1)+T(n-2),
\]
why can strong induction be useful?

\item State the Principle of Strong Induction.

\item Prove that every integer $n\geq2$ can be expressed as a product of primes.

\item Why can strong induction require multiple base cases?

\item What is the difference between ordinary and strong induction?
\end{enumerate}

\end{frame}

\begin{frame}{17. Strong Induction -- Answers}

\begin{enumerate}
\item The recurrence depends on multiple earlier values, so assuming all previous cases makes the induction step natural.

\item Assume $P(1),P(2),\ldots,P(k)$ are true and prove $P(k+1)$.

\item If $n$ is prime, it is already a product of one prime. If composite, write $n=ab$ where $a,b<n$ and apply the induction hypothesis to both.

\item The proof of $P(k+1)$ may depend on several earlier values rather than only $P(k)$.

\item Ordinary induction assumes only $P(k)$; strong induction assumes all previous cases up to $k$.
\end{enumerate}

\end{frame}

% =================================================
% 18. STRUCTURAL INDUCTION
% =================================================

\section{Structural Induction}

\begin{frame}{18. Structural Induction -- Questions}

\begin{enumerate}
\item \textbf{Computing:} For a binary expression tree, prove that the number of leaves equals the number of internal binary nodes plus one.

\item What is structural induction?

\item What are the main cases when proving a property of recursively defined binary trees?

\item How is structural induction different from ordinary induction?

\item Why is structural induction suitable for recursive data structures?
\end{enumerate}

\end{frame}

\begin{frame}{18. Structural Induction -- Answers}

\begin{enumerate}
\item Base case: a leaf has one leaf and zero internal nodes. For an internal node with two subtrees, combine the induction hypotheses:
\[
(L_1+L_2)=(I_1+1)+(I_2+1)
=I_1+I_2+2.
\]
Since the new internal node adds one,
\[
L=I+1.
\]

\item Structural induction proves properties by following the recursive construction of an object.

\item Base case: leaf. Recursive case: combine the two subtrees.

\item Ordinary induction is based on integer progression; structural induction is based on the construction rules of recursive objects.

\item Recursive structures such as trees, expressions, and lists are naturally generated by construction rules.
\end{enumerate}

\end{frame}

% =================================================
% FINAL REVISION SLIDE
% =================================================


% =================================================
% END
% =================================================

{
\setbeamertemplate{footline}{}
\begin{frame}
\centering
\vspace{1.5cm}

{\Huge\bfseries Thank You}

\vspace{0.5cm}

{\Large Mathematics of Computing}

\vspace{0.4cm}

{\large CUSAT M.Tech -- AI \& Software Engineering}

\vspace{0.8cm}

{\large Sayona C Sajan}

\end{frame}
}

\end{document}